% Source: https://artofproblemsolving.com/community/c3h1485435_funny_sequences
% Stable identity: aops-topic-1485435-q01
The idea is to group the terms in consecutive pairs; the pairs then form an ordinary arithmetic or geometric sequence.

\textbf{(a)} After each pair of steps the sequence increases by $1+2=3$. So the $k$th pair is $(3k-2,\ 3k-1)$: the terms are $a_{2k-1}=3k-2$ and $a_{2k}=3k-1$. In one formula, $a_n=\left\lfloor\dfrac{3n-1}{2}\right\rfloor$ (check: $n=1,2,3,4$ give $1,2,4,5$). The $100$th term is $a_{100}=3\cdot50-1=\boxed{149}$. Each pair sums to $6k-3$, so the first $100$ terms (that is, $50$ pairs) sum to
\[
\sum_{k=1}^{50}(6k-3)=6\cdot\dfrac{50\cdot51}{2}-150=7650-150=\boxed{7500}.
\]

\textbf{(b)} After each pair of steps the sequence is multiplied by $2\cdot4=8$. So the $k$th pair is $(8^{k-1},\ 2\cdot8^{k-1})$, with sum $3\cdot8^{k-1}$. The first $20$ terms are $10$ pairs, with sum
\[
3\left(1+8+\dots+8^9\right)=3\cdot\dfrac{8^{10}-1}{7}=\dfrac{3\cdot1073741823}{7}=\boxed{460175067}.
\]
